% Two-assembly architecture of the ejector unit operation (same
% ports-and-adapters split as the membrane unit, separate assemblies). The
% Python prototype is the reference implementation the core is ported from.
\begin{tikzpicture}[font=\small, >=Latex,
    box/.style={draw, thick, rounded corners, align=left, text width=5.0cm, inner sep=7pt}]
  \node[box, docnavy, fill=docnavy!6] (core)
    {\textbf{EjectorCore}\\[1pt]\footnotesize (physics, PME-agnostic)\\[3pt]
     \footnotesize $\bullet$ Expansion (nozzle, choke search)\\
     $\bullet$ Mixing (subsonic root), Diffuser\\
     $\bullet$ Rate / RateFlows (inverse rating)\\
     $\bullet$ Design / DesignForLift\\
     $\bullet$ Brent root and minimiser (SciPy ports)\\
     $\bullet$ IFlash port: PH, PS, PT, PQ at $\vec z$};
  \node[box, docred, fill=docred!4, right=2.6cm of core] (ad)
    {\textbf{Ejector.CapeOpen}\\[1pt]\footnotesize (COM / CAPE-OPEN adapter)\\[3pt]
     \footnotesize $\bullet$ EjectorUnitOperation\\
     $\bullet$ StreamAdapter (inlets; outlet PH flash)\\
     $\bullet$ PmeFlash : IFlash (scratch MO)\\
     $\bullet$ Real / Option / Array parameters\\
     $\bullet$ persistence, registration\\
     $\bullet$ HTML / text report};
  \node[box, docslate, fill=docslate!6, below=1.1cm of core, text width=5.0cm] (py)
    {\textbf{ejector1d.py}\\[1pt]\footnotesize reference implementation\\
     (Python + CoolProp): validation, calibration};
  \draw[->, thick, docgrey] (ad) -- (core) node[midway, above, font=\footnotesize] {uses};
  \draw[->, thick, docgrey, densely dashed] (py) -- (core)
    node[midway, right, font=\footnotesize] {ported to, regression-checked against};
  \node[font=\footnotesize, docgrey, below=2pt of ad] {net8.0-windows x64, COM via comhost};
\end{tikzpicture}
